\section{Negativity throughout a fully-entangled-fraction maximizing face}
\label{sec:negativity-face}

The saturation classification fixes the fully entangled fraction and the
$G$-concurrence, but it does not fix the negativity.  Here we solve the
minimum-negativity problem on every full-positive coherent-core rectangle,
and obtain both endpoints exactly on an infinite family of spectral
maximizer faces. We use the convention
$\mathcal N(\rho)=(\|\rho^\Gamma\|_1-1)/2$.
The conclusions concern the FEF-maximizing face under discussion; they do
not assert that its largest negativity is the largest among all states
with the same marginals.

Let $m_i>0$, $z=\sum_{i=1}^d\sqrt{m_i}\,|ii\rangle$, and let $I,J$ be
nonempty disjoint subsets of $\{1,\ldots,d\}$. Put
$L=\sum_i m_i$, $\delta=1-L>0$. Consider
\[
 \rho_R=|z\rangle\langle z|+R,
 \quad R\succeq0,\quad
 \operatorname{supp}R\subseteq\mathbb C^I\otimes\mathbb C^J,
 \quad \operatorname{Tr}_B R=\operatorname{diag}(u),\quad
 \operatorname{Tr}_A R=\operatorname{diag}(v),
\]
where $u_i,v_j>0$ and $\sum_i u_i=\sum_jv_j=\delta$.
This includes aligned fixed-matrix faces of full-rank spectral saturators,
as well as aligned fixed-population optimizer faces without any spectral
ordering assumption. No assertion is made that every prescribed-spectrum
optimizer admits this extra simultaneous alignment of its residual
marginals with the core. Without the spectral ordering condition, the
more general aligned face maximizes overlap with a fixed pure target;
FEF optimality over varying target bases is not asserted. Define
\[
 d_{ij}=\sqrt{m_i m_j},\qquad
 N_0=\sum_{1\le i<j\le d}\sqrt{m_i m_j}
                -\sum_{i\in I,j\in J}d_{ij}.
\]
The omitted pairs are distinct because $I\cap J=\varnothing$.

\subsection{A strict-convexity lemma for a saddle matrix}

\begin{lemma}\label{lem:strict-saddle}
For every positive definite $D\in M_n(\mathbb C)$, the function
\[
 C\longmapsto F_D(C)=
 \left\|\begin{pmatrix}C&D\\D&0\end{pmatrix}\right\|_1
\]
is strictly convex on the real vector space of Hermitian $n\times n$
matrices.
\end{lemma}
\begin{proof}
Write $H(C)$ for the displayed block matrix. It is invertible: the bottom
block of $H(C)(x,y)^T=0$ gives $Dx=0$, hence $x=0$ and then $Dy=0$.
For an invertible Hermitian matrix $H$, the maximizer of
$\operatorname{Tr}(SH)$ over Hermitian contractions $-I\preceq S\preceq I$
is unique and equals $\operatorname{sign}(H)$. This follows by diagonalizing
$H$: each diagonal entry of a maximizing contraction must have the sign
of the corresponding eigenvalue, and a contraction with a diagonal entry
of modulus one has zero off-diagonal entries in that row and column.

Suppose equality holds in convexity between $H(C_1)$ and $H(C_2)$ at an
interior mixing parameter. The unique supporting contraction of their
invertible mixture maximizes both endpoints, so
$S=\operatorname{sign}H(C_1)=\operatorname{sign}H(C_2)$.
Write $S=\begin{psmallmatrix}A&B\\B^*&E\end{psmallmatrix}$.
Since $H(C_1)S=|H(C_1)|\succ0$, its lower-right block $DB$ is positive
definite. Thus $B$ is invertible. Both $H(C_i)$ commute with $S$, hence
$\operatorname{diag}(C_1-C_2,0)$ commutes with $S$. Its upper-right
commutator block is $(C_1-C_2)B=0$. Invertibility of $B$ gives
$C_1=C_2$, proving strict convexity.
\end{proof}

Positive definiteness of $D$ matters. For
$D=\operatorname{diag}(1,0)$ and $C_t=\operatorname{diag}(0,t)$,
$F_D(C_t)=2+t$ for $t\ge0$; strict convexity fails.

\subsection{The unique minimum is a classical transportation problem}

\begin{theorem}\label{thm:negativity-transport}
The map $R\mapsto\mathcal N(\rho_R)$ is strictly convex. Its unique minimum
on the prescribed residual-marginal face is attained by a diagonal
residual
\[
 R_* = \sum_{i\in I,j\in J}t^*_{ij}|ij\rangle\langle ij|,
\]
where $t^*$ is the unique solution of the strictly convex transportation
problem
\[
 \boxed{\quad
 \min_{\substack{t_{ij}\ge0,\ \sum_jt_{ij}=u_i\\
                        \sum_it_{ij}=v_j}}
 \left[N_0+\frac12\sum_{i\in I,j\in J}
       \left(\sqrt{t_{ij}^2+4m_i m_j}-t_{ij}\right)\right].
 \quad}
\]
For every nondiagonal $R$, deleting its off-diagonal entries strictly
reduces $\mathcal N(\rho_R)$ and preserves both actual marginal matrices.
\end{theorem}
\begin{proof}
In the ordered bases of $I\times J$ and the reversed rectangle
$J\times I$, the varying partial-transpose block is
\[
 H(C)=\begin{pmatrix}C&D\\D&0\end{pmatrix},
 \quad C=R^\Gamma,\quad D=\operatorname{diag}(d_{ij})\succ0.
\]
All other diagonal and two-dimensional blocks are fixed. Consequently
\[
 \mathcal N(\rho_R)=N_0+\frac{F_D(R^\Gamma)-\delta}{2}.
\]
Partial transpose is an injective real-linear map on Hermitian matrices,
so Lemma~\ref{lem:strict-saddle} implies strict convexity.

Average the local unitaries $U\otimes\overline U$ over all diagonal
sign matrices $U$. They fix $z$, preserve the diagonal marginal matrices,
and delete precisely the off-diagonal entries of $R$. The characters
$\varepsilon_i\varepsilon_j$ on the disjoint rectangle are distinct.
Negativity is invariant under each local unitary and convex under their
mixture. Strict convexity makes the inequality strict unless every
conjugate is the same, which is equivalent to $R$ being diagonal.
For a diagonal residual the paired block eigenvalues are
$(t_{ij}\pm\sqrt{t_{ij}^2+4d_{ij}^2})/2$; summing their negative parts gives
the displayed objective. Each scalar cost has second derivative
$2d_{ij}^2/(t^2+4d_{ij}^2)^{3/2}>0$.
The compact nonempty transportation polytope therefore has exactly one
minimizer, which is also the unique quantum minimizer.
\end{proof}

This gives an exact efficient convex characterization even when the
minimizing transport is not a product coupling. Its optimality conditions
are particularly simple: there exist real potentials $\lambda_i,\mu_j$
such that, with
$g_{ij}'(t)=\tfrac12(t/\sqrt{t^2+4d_{ij}^2}-1)$,
\[
 \lambda_i+\mu_j=g'_{ij}(t^*_{ij})\quad(t^*_{ij}>0),
 \qquad \lambda_i+\mu_j\le-\tfrac12\quad(t^*_{ij}=0).
\]
These conditions are necessary and sufficient. For positive cells set
$h_{ij}=1+2\lambda_i+2\mu_j\in(0,1)$; then
$t^*_{ij}=2d_{ij}h_{ij}/\sqrt{1-h_{ij}^2}$.

\begin{corollary}[A closed sharp lower bound]\label{cor:negativity-total}
Put $P=(\sum_{i\in I}\sqrt{m_i})(\sum_{j\in J}\sqrt{m_j})$.
Then every residual satisfies
\[
 \boxed{\quad \mathcal N(\rho_R)\ge
 N_0+\frac{\sqrt{\delta^2+4P^2}-\delta}{2}.\quad}
\]
Equality holds exactly when
\[
 R=\sum_{i,j}\frac{\delta\sqrt{m_i m_j}}{P}
                         |ij\rangle\langle ij|.
\]
In particular, equality is feasible exactly when
$u_i=\delta\sqrt{m_i}/\sum_{k\in I}\sqrt{m_k}$ and
$v_j=\delta\sqrt{m_j}/\sum_{\ell\in J}\sqrt{m_\ell}$.
\end{corollary}
\begin{proof}
First dephase as in the theorem. The triangle inequality in $\mathbb R^2$
gives
$\sum_{i,j}\sqrt{t_{ij}^2+4d_{ij}^2}\ge
\sqrt{(\sum t_{ij})^2+4(\sum d_{ij})^2}$.
All $d_{ij}>0$, so equality holds exactly when $t_{ij}/d_{ij}$ is constant.
Trace $\delta$ fixes that constant to $\delta/P$. Strictness under
dephasing gives the full quantum equality statement.
\end{proof}

\paragraph{An exact nonproduct minimizer.}
The usual product residual need not minimize negativity, even within a
spectral FEF-maximizer face. Put
\[
 k=\frac{17}{3}+\frac4{\sqrt{21}},\quad
 a=\left(\frac{101}{50}+\frac{k}{10}\right)^{-1},\quad
 b=\frac a{100},\quad d_0=\frac a{10},\quad \delta=kd_0.
\]
Take $m=(a,a,b,b)$, $I=\{1,2\}$, $J=\{3,4\}$, and let
\[
 t^*=d_0\begin{pmatrix}4/\sqrt{21}&3/2\\3/2&8/3\end{pmatrix},
 \qquad u=v=t^*\binom11.
\]
Normalization follows from $2a+2b+\delta=1$. Both prescribed marginals
can be put in decreasing order by permutations within the two equal-core
blocks. Indeed $b+\max_jv_j=(1/100+5/12)a<a$.
The KKT potentials $\lambda=(-1/5,-1/10)$ and $\mu=(-1/10,0)$ have
sums $(-3/10,-1/5;-1/5,-1/10)$, exactly the four cost derivatives at
$t^*$. Thus the unique minimum is
\[
 \mathcal N_{\min}=a+b+d_0\left(\frac43+\frac3{\sqrt{21}}\right).
\]
But
$\det(t^*/d_0)=32/(3\sqrt{21})-9/4>0$, since
$128^2>27^2\cdot21$. The product transport $uv^T/\delta$ has determinant
zero and is different. Strict convexity therefore makes its negativity
strictly larger. This example distinguishes the coherent-minimum product
construction, which proves sharpness of the FEF bound, from the residual
that minimizes a different entanglement measure on the same optimal face.

\subsection{A quantitative penalty for residual coherence}

There is also an explicit, nonoptimal stability estimate. Let
$d_- =\min_{i,j}d_{ij}>0$, $d_+=\max_{i,j}d_{ij}$ and
\[
 h=\frac{\sqrt{\delta^2+4d_-^2}-\delta}{2},\qquad
 \kappa=\frac{h^2}{8d_+^2(\delta+d_+)}.
\]
For the diagonal dephasing $\Delta R$,
\begin{equation}\label{eq:negativity-coherence-gap}
 \mathcal N(\rho_R)-\mathcal N(\rho_{\Delta R})
       \ge \kappa\,\|R-\Delta R\|_2^2,
\end{equation}
where $\|\cdot\|_2$ is the Hilbert--Schmidt norm.
To prove this, along every residual segment we have
$\|C\|\le\|C\|_2=\|R\|_2\le\delta$ and
$\|H(C)\|\le M=\delta+d_+$.
For an eigenpair $H(C)(x,y)=\lambda(x,y)$, $x\ne0$ and
\[
 \lambda^2\|x\|^2-\lambda\langle x,Cx\rangle-\|Dx\|^2=0.
\]
It follows that $|\lambda|^2+\delta|\lambda|\ge d_-^2$, hence
$|H(C)|\succeq hI$. If
$S=\operatorname{sign}H=\begin{psmallmatrix}A&B\\B^*&E\end{psmallmatrix}$,
then $DB=(|H|)_{22}\succeq hI$, so
$\sigma_{\min}(B)\ge h/d_+$.
For a Hermitian perturbation $Q$ in the upper-left block, the spectral
second-derivative formula gives
\[
 F_D''(C)[Q,Q]
 =4\sum_{\lambda_r>0>\lambda_s}
      \frac{|\langle r,\operatorname{diag}(Q,0)s\rangle|^2}
           {\lambda_r-\lambda_s}
 \ge\frac{\|[\operatorname{diag}(Q,0),S]\|_2^2}{4M}
 \ge\frac{h^2}{2Md_+^2}\|Q\|_2^2.
\]
The same formula is valid with repeated eigenvalues, because the positive
and negative spectral subspaces stay separated from zero; repeated values
within one sign contribute zero. Averaging this strong-convexity inequality
over the finite diagonal-sign twirl gives
$F_D(C)-F_D(\Delta C)\ge h^2\|C-\Delta C\|_2^2/(4Md_+^2)$.
Negativity is one half of this varying trace norm and partial transpose
preserves the Hilbert--Schmidt norm, proving
\eqref{eq:negativity-coherence-gap}.

\subsection{An exact negativity interval in every even dimension}

\begin{theorem}\label{thm:balanced-negativity}
Fix $q\ge2$, $a,b>0$ and
$\delta=1-q(a+b)>0$ with $\delta\le q(a-b)$.
On $\mathbb C^{2q}\otimes\mathbb C^{2q}$ set
\[
 z=\sqrt a\sum_{i=1}^q|ii\rangle+
   \sqrt b\sum_{j=q+1}^{2q}|jj\rangle,
 \qquad \rho_\omega=|z\rangle\langle z|+\delta V\omega V^*,
\]
where $V|ij\rangle=|i,j+q\rangle$ and
$\omega\succeq0$ has both marginals $I_q/q$.
All these states have the same actual marginal matrices, decreasing spectra
\[
 \alpha=((a+\delta/q)^{[q]},b^{[q]}),\quad
 \beta=(a^{[q]},(b+\delta/q)^{[q]}),
\]
and attain the same sharp spectral FEF bound
\[
 f(\rho_\omega)=\frac q2(\sqrt a+\sqrt b)^2.
\]
Their negativity takes every value in the exact interval
\[
 \boxed{\quad N_{\min}\le\mathcal N(\rho_\omega)\le N_{\max},\quad}
\]
where
\begin{align*}
 N_{\min}&=\frac{q(q-1)}2(a+b)
          +\frac{\sqrt{\delta^2+4q^4ab}-\delta}{2},\\
 N_{\max}&=\frac{q(q-1)}2(a+b)
          +\frac{q\sqrt{\delta^2+4q^2ab}-\delta}{2}.
\end{align*}
The lower endpoint is attained only by $\omega=I_{q^2}/q^2$.
The upper endpoint is attained exactly by pure maximally entangled
$\omega$. Thus the two endpoints have joint-state ranks $1+q^2$ and $2$,
respectively, despite having identical marginals, FEF, Schmidt number
$2q$, $G$-concurrence $2q\sqrt{ab}$, and partial-transpose inertia.
\end{theorem}
\begin{proof}
The ordering assumption is equivalent to $a\ge b+\delta/q$ and ensures
that the two decreasing marginal lists have componentwise minimum
$m=(a^{[q]},b^{[q]})$. The saturation normal form gives the FEF formula
and the common entanglement invariants. The varying block has
$D=\sqrt{ab}I_{q^2}$ and $C=\delta\omega^\Gamma$.
Diagonalizing $C$ splits it into two-dimensional blocks, giving
\[
 \mathcal N(\rho_\omega)=\frac{q(q-1)}2(a+b)
   +\frac12\left(\operatorname{Tr}\sqrt{C^2+4abI}-\delta\right).
\]
Strict scalar convexity of $s\mapsto\sqrt{s^2+4ab}$, and
$\operatorname{Tr}C=\delta$, give the lower endpoint, with equality only
for $C=\delta I/q^2$.
For the upper endpoint use Cauchy--Schwarz and preservation of the
Hilbert--Schmidt norm by partial transpose:
\[
 \operatorname{Tr}\sqrt{C^2+4abI}
 \le \sqrt{q^2(\operatorname{Tr}C^2+4abq^2)}
 \le q\sqrt{\delta^2+4abq^2}.
\]
The second inequality is $\operatorname{Tr}\omega^2\le1$. Equality forces
$\omega$ to be pure. Its fixed uniform marginals then force maximal
entanglement, and its partial transpose has squared value $I/q^2$, so
both inequalities are equalities. This also proves the complete upper
endpoint classification. Finally, the path
$\omega_t=(1-t)I/q^2+t|\Phi_q\rangle\langle\Phi_q|$, $0\le t\le1$,
is feasible and continuous and joins the endpoints. Thus every
intermediate negativity is attained.
\end{proof}

Along that path a fully explicit formula is available. Put
$n_\pm=q(q\pm1)/2$ and
$c_+(t)=\delta[1+(q-1)t]/q^2$,
$c_-(t)=\delta[1-(q+1)t]/q^2$. Then
\[
 \mathcal N(\rho_{\omega_t})=
 \frac{q(q-1)}2(a+b)+\frac12\left[
 n_+\sqrt{c_+(t)^2+4ab}+n_-\sqrt{c_-(t)^2+4ab}-\delta\right].
\]
It is strictly increasing for $0<t\le1$: strict convexity gives a strictly
increasing derivative and the derivative vanishes at $t=0$.

\paragraph{Exact four-level example.}
Take $q=2$, $a=3/10$, $b=1/10$, $\delta=1/5$. Then
\[
 \alpha=(2/5,2/5,1/10,1/10),\quad
 \beta=(3/10,3/10,1/5,1/5),\quad
 f=\frac25+\frac{\sqrt3}{5},\quad G_4=\frac{2\sqrt3}{5},
\]
while
\[
 \boxed{\quad 1\le\mathcal N\le\frac3{10}+\frac{\sqrt{13}}5.\quad}
\]
The lower state has rank five and a classical uniform residual; the upper
states have rank two and a Bell residual. This is an exact counterexample
to the assertion that the FEF, both full marginal spectra, the
$G$-concurrence, and partial-transpose inertia determine the negativity.

\paragraph{Scope and prior art.}
The negativity and its local-unitary invariance and convexity are standard;
see Vidal--Werner, \emph{A computable measure of entanglement},
\href{https://arxiv.org/abs/quant-ph/0102117}{arXiv:quant-ph/0102117}.
Fixed-marginal negativity optimization has a substantial prior literature,
including Baio et al., \emph{Bounds on bipartite entanglement from fixed
marginals}, \href{https://arxiv.org/abs/1904.07650}{arXiv:1904.07650}.
The new candidates here are the strict-convexity reduction on an entire
FEF-maximizing face, its complete quantum minimum and quantitative
coherence penalty, and the exact all-even-dimensional endpoint family.
A bounded primary-source search did not locate these precise statements;
this is not a certified worldwide-priority claim. These are exact written
proofs; any accompanying scalar Lean certificates must not be described
as an end-to-end formalization of partial transpose and trace norm.
